% Generated from the Notion page by tools/notion2tex.py (see report/README.md).
\section{Discussion}\label{sec:6}

\subsection{Results against the goals}\label{sec:6.1}

Table~\ref{tab:21} closes the loop from the goals of Section~\ref{sec:1} to the evidence of Section~\ref{sec:5}. Every row of the specification is met except that G2 has no numeric target by construction; its row reports the achieved value and the level of evidence at which it was obtained.

\begin{table}[!htbp]
\centering
\footnotesize
\caption{Goals of Section~\ref{sec:1} traced to the specification used in the search (Section~\ref{sec:3}, Table~\ref{tab:2}) and to the evidence that closes each one.}
\label{tab:21}
\begin{tabularx}{\linewidth}{L{2.1cm}L{3.4cm}L{2.0cm}YL{2.0cm}}
\toprule
Goal & Specification & Evidence level & Result & Where \\
\midrule
G1 interface and network & \texttt{srcnn} signature with float arrays, parameters read from DDR on every call, 1-64-32-1, replicate padding & Board & Parameter-reload test returns the analytically expected 104/256; original parameters restore the bit-exact output & Section~\ref{sec:5}, Table~\ref{tab:12} \\
G1 correctness & Bit-exact against the fixed-point reference on all 13 Set14 images & Board & 13 / 13 bit-exact, plus Butterfly (Set5) & Table~\ref{tab:13} \\
G1 quality & Mean PSNR on Set14 no worse than $-$0.1 dB against float & Board & 29.083 dB on the board, identical to co-simulation; $-$0.076 dB against the float golden; worst image $-$0.268 dB (foreman) & Table~\ref{tab:8}, Table~\ref{tab:12} \\
G2 speed & ms per frame, as low as possible & RTL $\times$ routed clock & 1.365 ms (273,023 cycles at 5.0 ns) & Section~\ref{sec:3}, Section~\ref{sec:4} \\
G2 speed & same & Board, hardware counter & 2.938 ms (489,585 cycles at 166.665 MHz), 762$\times$ faster than the 2.238 s baseline measured on the same board & Table~\ref{tab:19}, Table~\ref{tab:20} \\
G3 fit & Post-route timing met, zero unrouted nets, resources within the device, as an integrated system & Vivado implementation of the PS+PL system & Met at 5.333 ns (WNS +0.043 ns); 1,032 DSP (82.7 \%), 63,078 LUT (53.9 \%), 54.5 BRAM36 (37.9 \%), 97.7 \% of CLBs occupied & Table~\ref{tab:11} \\
G4 evidence & Every number traceable to an archived run with its evidence level & Repository & 231 registered experiments, 71 design directories, registry generated from reports, evolution graph & Section~\ref{sec:3}, Figure~\ref{fig:2} \\
\bottomrule
\end{tabularx}
\end{table}

Two things in the table deserve comment. The board figure for G2 is 2.15$\times$ the RTL figure; Section~\ref{sec:6.3} shows that this is not a defect of the datapath but of the memory path and the clock configuration, and both are measured rather than assumed. And the G3 row shows a design that fits the device but only just: the DSP budget is 83 \% used and the logic is spread over 98 \% of the CLBs, which is the reason the final gate of the search had to be a routed result and not an HLS estimate.

\subsection{Where the speedup comes from}\label{sec:6.2}

The 762$\times$ between the baseline and the final design is not one idea. It is a chain of structural changes, each of which was measured against a control at the time (Section~\ref{sec:3}), and the chain can be read off the milestones of the search. Table~\ref{tab:22} lists them with the level of evidence of each number; the HLS column and the routed column are kept apart because they differ by about two times for the same design.

\begin{table}[!htbp]
\centering
\footnotesize
\caption{Milestones from the original source to the board. ``HLS'' is estimated cycles $\times$ target clock; ``RTL $\times$ routed'' is co-simulated cycles $\times$ an independently closed post-route clock; ``Board'' is the hardware cycle counter at the configured PL clock. Rows 1 and 3 to 9 are from the experiment registry; row 2 is the hand-written baseline of Section~\ref{sec:2}.}
\label{tab:22}
\begin{tabularx}{\linewidth}{rL{3.0cm}YL{2.3cm}R{1.3cm}L{1.7cm}}
\toprule
\# & Design & What changed relative to the previous row & Evidence & ms / frame & DSP \\
\midrule
1 & E000, original float source, 10 ns target & --- & HLS & 19,601 & --- (11,064 BRAM18K, does not fit) \\
2 & Baseline (Section~\ref{sec:2}): 16$\times$16 tiles, channel parallelism, pipelining, float & Tiling makes it fit; layers still sequential & Board & 2,238 & 46 (3.69 \%) \\
3 & E039: streaming dataflow, 10-bit fixed point, T = 32, 10 ns & Whole-network dataflow, no intermediate maps, one DSP per product & HLS & 29.83 & 251 \\
4 & E043: T = 16, 3.33 ns target & Fold halved, clock target tightened & HLS & 6.94 & 505 \\
5 & F25 / F92: continuous direct convolution, loader off the critical path, T = 16 & Window multiplexers removed; loader restructured & HLS / RTL $\times$ routed at 5.0 ns & 2.64 / 5.34 & 592 \\
6 & W21 / W32: Winograd on Conv1, T = 8 & Conv1 products 5,184 $\rightarrow$ 2,304 per pixel, so T = 8 fits & HLS / RTL $\times$ routed at 4.5 ns & 1.75 / 2.51 & 644 \\
7 & C16 / C21: window FIFO 16 deep, free-running pipelines, T = 8 & Stalls on the transformed-window FIFO removed & RTL $\times$ routed at 4.5 ns & 2.41 & 644 \\
8 & C13 / C18: T = 4 with 8-bit packed weights on Conv2 & Two Conv2 products per DSP, so T = 4 fits in 1,032 DSPs & HLS / RTL $\times$ routed at 5.0 ns & 1.05 / 1.50 & 1,032 \\
9 & \textbf{C17 / C22: window FIFO 32 deep (selected)} & One change: FIFO depth 2 $\rightarrow$ 32 & HLS / RTL $\times$ routed at 5.0 ns & 1.05 / \textbf{1.365} & 1,032 \\
10 & C22 on the board, first bitstream & PL clock actually 125 MHz, three AXI masters on HP0 & Board & 5.05 & 1,032 \\
11 & C22, PL0 divider 8 $\rightarrow$ 6 & 166.665 MHz & Board & 3.984 & 1,032 \\
12 & \textbf{C22, output on its own HP port} & Write bursts no longer contend with reads & Board & \textbf{2.938} & 1,032 \\
\bottomrule
\end{tabularx}
\end{table}

Read as three blocks, the chain is:

\textbf{From the baseline to the first routed streaming design, 419$\times$ (rows 2 to 5).} This is the largest block and it comes from changing the shape of the computation, not from tuning it. The baseline runs the three layers in sequence inside each 16$\times$16 tile, re-fetches a 28$\times$28 window from DDR per tile, keeps 46 DSPs busy and carries floating-point arithmetic, whose multiply-add on this device needs several DSPs and a deep pipeline per operation. The streaming design replaces all of that: the three layers run concurrently on different pixels of the same frame, the 23.81 MiB of intermediate feature maps that the direct form needs (and that E020 to E032 showed cannot be pragma'd away: 11,064 BRAM18K regardless of directives) are reduced to a few line buffers and FIFOs, and 10-bit fixed point lets one DSP48E2 compute one product per cycle so that 592 of them can work at once. The precision was chosen before the architecture and then frozen (Table~\ref{tab:8}): 10-bit weights and activations cost 0.070 dB of mean PSNR, uniform 8-bit would have cost 5.18 dB, and every later design had to match the fixed-point reference bit for bit, so that the whole chain below is attributable to structure and not to numerics. Within this block the HLS numbers show the order of the gains even though they are not board-comparable: streaming (rows 1 to 3), then halving the fold and tightening the clock (row 4, 4.3$\times$), then the continuous engine and the loader (row 5, 2.6$\times$ at the HLS level). The loader step is worth singling out because its evidence was a resource breakdown rather than a timing number: in F25, 46 \% of the LUTs sat outside the dataflow region in the multiplexers between the weight banks and the engines, and streaming the parameters instead (F27) cut those out-of-region LUTs from 32 k to 3.2 k, which is what later made the larger designs routable.

\textbf{From the first routed design to the selected one, 3.9$\times$ at the same clock (rows 5 to 9).} The routed clock is 5.0 ns in both rows 5 and 9, so this whole factor is cycles: 1,068,183 $\rightarrow$ 273,023, which is almost exactly the fold going from T = 16 to T = 4. The fold is bounded by DSPs, and Table~\ref{tab:5} shows that direct convolution cannot go below T = 8 on this device (T = 8 direct needs 1,004 DSPs, F29 and F30; T = 4 direct needs 2,008). The two arithmetic reductions are what make T = 4 reachable, and the controls say exactly what each one does and does not do:

\begin{itemize}
\item \emph{Winograd is an enabler, not a speedup by itself.} At T = 16 the Winograd network W18 uses 326 DSPs instead of 502 but takes the same number of cycles as the direct design (1,056,850 against 1,048,687), because the frame time is set by the fold, not by the multiplier count. The processing-element pair W01 and W02 measured the trade at the smallest scale: 16 DSPs instead of 36 per 2$\times$2 tile, for roughly 1,000 extra LUTs of transform adders (1,499 against 503). Winograd pays off only where the direct design no longer fits, which is at T = 8 (644 DSPs against 1,004) and T = 4 (576 for Conv1 against 1,296). The 2.25$\times$ reduction in multiplications is the textbook figure for F(2$\times$2, 3$\times$3) \cite{ref16}; the design keeps it exact in integers by scaling the kernel transform by two and shifting the result, which is why the bit-exact gate holds through this layer.
\item \emph{Packing is what closes the last gap.} With Winograd on Conv1 and direct Conv2, T = 4 still needs 1,288 DSPs, 40 more than the device has. Packing two Conv2 weights into one 27-bit DSP operand halves that layer's DSPs to 256 and brings the total to 1,032. The exhaustive processing-element check P01 and the mixed-precision runs E062 and E065 bound its cost: it is exact for signed 8-bit weights, and 8-bit weights on Conv2 alone cost 0.006 dB more than the 10-bit baseline, against the 5 dB that uniform 8-bit would have cost. The alternative that keeps 10-bit weights everywhere (P07's residual correction) is also exact but costs 454 LUTs per product pair instead of 132, and applied to the whole network at T = 4 (P12) it needs 292 k LUTs, 2.5$\times$ the device.
\item \emph{Fold balance.} Once the three engines are at T = 4, their iteration counts are within 1 \% of each other (262,144, 260,100, 260,100), and the measured frame (273,023 cycles) is within 4 \% of the slowest engine's bound. The RTL design is therefore compute-bound and balanced; there is no further cycle to be had at T = 4 without changing the fold, and T = 2 is out of reach by a factor of two on DSPs with any combination of techniques (Table~\ref{tab:5}).
\item \emph{FIFO depth.} The remaining 10 \% between rows 8 and 9 came from one parameter. Co-simulation of W21 (F97) had shown 45,974 stall cycles on the transformed-window FIFO; deepening it (F47 at T = 8, then C17 at T = 4) took the T = 4 design from 300,632 to 273,023 cycles. Only co-simulation sees this: HLS estimates the same cycle count for C13 and C17.
\end{itemize}

What did \emph{not} move in this block is the clock. The routed period stayed at 5.0 ns from F92 to C22 while HLS estimated 2.4 to 2.9 ns for the same designs, and the three attempts to close 4.5 ns on the T = 4 RTL (W36, C36, C40) all failed timing with their worst paths in the Winograd arithmetic. The calibration data behind the gate change of Section~\ref{sec:3} (routed LUTs 0.46 to 0.66 of the HLS estimate, routed clock 1.7 to 2$\times$ the estimate across 18 implemented designs) is the quantitative form of that observation.

\textbf{From the RTL to the board, $\div$2.15 (rows 9 to 12).} The last block is a loss, and Section~\ref{sec:6.3} decomposes it: 0.27 ms is the PL clock at 166.665 MHz instead of 200 MHz, and 1.30 ms is single-beat AXI reads on the input and parameter paths. Of the 5.05 ms of the first bitstream, 1.07 ms was recovered by reading the actual IOPLL configuration and changing one divider, and 1.05 ms by giving the output master its own HP port so that its writes no longer contend with the two read masters \cite{ref20}. Multiplying the three blocks together, 419 $\times$ 3.9 $\div$ 2.15, gives the 762$\times$ of Table~\ref{tab:19}.

\textbf{What the final number means.} On the board the accelerator performs the 522 million MACs of the original network in 2.938 ms, an effective 178 GMAC/s, against a raw capacity of 1,032 DSPs $\times$ 166.665 MHz = 172 G multiplications per second; the effective rate exceeds the raw one because Winograd and packing remove multiplications rather than perform them. The ratio of RTL to on-board cycles, 273,023 / 489,585 = 56 \%, is the fraction of board cycles in which the datapath does the work it was built for. At the RTL level the same datapath is at 96 \% of its structural bound. The design is therefore finished as a datapath and unfinished as a system, and that is the order in which the next iteration should spend its effort.

\subsection{Where the time goes between the RTL and the board}\label{sec:6.3}

The last block of Table~\ref{tab:22} is the only one in which the design got slower, and because it was invisible at every evidence level available during the search, it is worth decomposing in full.

The final architecture needs 273,023 cycles in RTL co-simulation but 489,585 cycles on the board, and the measurements of Section~\ref{sec:5} allow the difference between the 1.365 ms RTL figure and the 2.938 ms hardware time to be decomposed into two parts with different causes.

\textbf{The clock accounts for 0.27 ms.} The RTL figure assumes the routed 200 MHz clock; at the 166.665 MHz the board runs at, the same 273,023 cycles take 1.638 ms. It would close completely at a divider of 5 if the integrated system meets timing at 200 MHz (Section~\ref{sec:6.5}).

\textbf{The memory path accounts for the remaining 1.30 ms, or 216,562 cycles.} First, the clock sweep (Table~\ref{tab:15}) fits \emph{t = a + b/f} with \emph{a} $\approx$ 1.0 ms and \emph{b} $\approx$ 493 k cycles in the HP0 configuration: about 1 ms of the frame did not scale with the PL clock. After the output was moved to HP1, the hardware time of 489,585 cycles $\times$ 6 ns = 2.937 ms coincides with the \emph{b} term alone, so the split removed almost exactly the clock-independent component. The natural reading, consistent with the AMD note on HP port arbitration, is that the fixed \textasciitilde{}1 ms was the output master's write bursts contending with the two read masters for one HP port and one DDR arbitration slot; with the writes on their own port, this contention disappeared and the 26.6 \% reduction of Table~\ref{tab:18} followed. Second, what remains scales with the PL clock, and the read-path statistics (Figure~\ref{fig:10}) locate it: over 77.7 \% of the frame the image channel had a read outstanding with no data returned, and both read channels averaged one beat per burst. Third, the arithmetic is consistent: the input image is 255 $\times$ 255 = 65,025 float words and the parameters are 8,129 words (5,184 + 64, 2,048 + 32, 800 + 1), so the frame issues about 73,150 single-beat reads, and 216,562 extra cycles spread over them is an average of 3.0 stall cycles per read. That is exactly the signature of a loader that pays the full interconnect and DDR latency on every word and overlaps only a few transactions, instead of amortising one address phase over a long burst.

\textbf{Why the stalls reach the engines.} In co-simulation (Figure~\ref{fig:4}) the pipeline is compute-bound: the engines wait less than 5 \% of the time while the producers in front of them are write-blocked 73 \% of the time. The input loader is the first of those producers, and on the board its reads do not arrive fast enough to keep its FIFO full. The condition inverts: the loader becomes read-blocked, the window generator and the Winograd engine starve, and because all three layers are coupled by finite FIFOs the stall propagates through the whole dataflow region. A memory-side stall therefore costs frame time one-for-one, which is why 216 k read-stall cycles appear as 216 k extra frame cycles rather than being hidden behind compute.

\textbf{Why HLS did not produce bursts.} The \texttt{max\_read\_burst\_length=256} setting only caps the burst that HLS may generate; Vitis HLS infers a burst only when it can prove that the accesses to an \texttt{m\_axi} port inside a loop are consecutive and that nothing in the loop body can stall between them. The C22 loader is three lines (Listing~\ref{lst:2}): a single pipelined loop over the 65,025 pixels that reads \texttt{in[0][i /\allowbreak{} W][i \% W]}, encodes the float to 10-bit fixed point and writes the result to the output stream in the same iteration. Two things defeat burst inference here. The address is formed with a division and a modulo by W = 255, which is sequential in fact but not provably so to the compiler, and the external read shares its loop body with a blocking \texttt{stream.write}, so the read of pixel \emph{i} + 1 cannot be issued until the write of pixel \emph{i} has been accepted. HLS therefore generates one AXI transaction per word, with ARLEN = 0, which is exactly what the counters report. The parameter loaders have the same structure.

\begin{lstlisting}[language=C++,caption={The input loader of C22 (\texttt{read\_input} in the original source): one pipelined loop that reads a word from DDR, encodes it and writes it to the stream in the same iteration.},label={lst:2}]
static void read_input(ftmap_t in[N0][H][W], hls::stream<input> &stream) {
    read_pixels: for (int i = 0; i < H * W; ++i) {
#pragma HLS PIPELINE II=1
        stream.write(encode_input(in[0][i / W][i % W]));
    }
}
\end{lstlisting} A trial restructuring (\texttt{burst\_\allowbreak{}input\_\allowbreak{}trial/\allowbreak{}}) that splits the loop into a load-only row loop into a 255-word BRAM buffer followed by a separate encode loop confirms the diagnosis: with no other change, HLS reports \texttt{HLS 214-115: multiple burst reads of length 65025} on \texttt{image\_in}, C simulation remains bit-exact, and the synthesis estimate is unchanged (263,092 cycles, 1,032 DSPs, +1 BRAM18K). This variant has not yet been co-simulated or run on the board, so its effect on the frame time is a prediction, not a result.

Co-simulation cannot reveal this, because its AXI model answers every read after a small fixed latency with no arbitration; the design is in fact bandwidth-starved only on real silicon. This is the concrete instance of the limitation noted in Section~\ref{sec:3}: during the search, every evidence level up to post-route implementation was available, and none of them could see the 2.2$\times$ that separates the RTL figure from the board.

\textbf{What the number means.} Measured on the board, the accelerator completes the 522 million multiply-accumulates of the original network in 2.938 ms, an effective rate of 178 GMAC/s, against a raw DSP capacity of 1,032 $\times$ 166.665 MHz = 172 G multiplications per second; the effective rate exceeds the raw one because Winograd and DSP packing remove multiplications rather than perform them. The RTL cycle count divided by the on-board count, 273,023 / 489,585 = 56 \%, is the fraction of on-board cycles in which the datapath is doing the work it was designed for; the remaining 44 \% is the price of the memory path, and it is the obvious target for the next iteration. Against the baseline, the speedup of 762$\times$ comes from three multiplicative sources: fixed-point and whole-network dataflow (the layers no longer run sequentially and nothing leaves the chip between them), Winograd and packing (fewer DSPs per product, so T = 4 fits in 1,032 DSPs where direct convolution would not fit at all), and the clock (166.665 MHz against the baseline's lower clock). The DSP utilisation rose from 3.69 \% to 82.7 \% in the same step.

The conclusion for the method is that the computational architecture obtained through HLS optimisation cannot determine the final performance alone; the memory access pattern and the system interconnect in the actual deployment are equally important, and they must be measured on the board because no earlier level of evidence shows them.

\subsection{On the method}\label{sec:6.4}

Section~\ref{sec:3} described the protocol and what it caught. This section asks whether it was worth it, what it did not do, and how far the findings generalise.

\textbf{What the protocol bought.} Three things that a hand search of the same length would probably not have produced. The first is the calibration of HLS estimates against routed results. It exists only because every run was archived with its evidence level and because the rule ``a lower level may not be reported as a higher one'' forced the HLS and routed numbers for the same design to sit side by side; the moment five such pairs existed the 2$\times$ pattern was obvious, and acting on it (the gate change of 2026-10-12) is what admitted the T = 4 designs that an HLS-ranked search would have discarded on LUT count (C13: 127,966 LUTs estimated, 58,529 routed). The second is the negative results. Of 231 registered experiments, 100 ended in failure, infeasibility, invalid configuration or timeout, and several of those closed a direction for good: output-tile fusion (S03 to S06, 130 to 228 ms), DDR scratch buffers (S01, whose ports were never exported and silently became 11,286 BRAM18K), scatter-based Conv3 (G03, 1.56$\times$ the cycles of the gather control), and the naive combination of Winograd with packing (C11, 249 k LUTs). Each is recorded with its card, its prediction and its measurement, so the report can say why the final design has the shape it has and not only that it works. The third is attribution: because every card changes one thing against a named control, Table~\ref{tab:22} and the bullets of Section~\ref{sec:6.2} can be written at all.

\textbf{What it cost.} The search ran for about four days of wall-clock time on a 12-core, 15 GiB host that was busy essentially continuously, with timeouts of 60 minutes for synthesis, 6 hours for co-simulation and 3 hours for implementation; the archive is several gigabytes; and the humans spent roughly one review per 10 to 20 experiments reading interpretations and rewriting plans. Against that, the hand-written baseline took the same two people from 2026-09-29 to the start of the search to reach 2.238 s, a 30$\times$ cycle reduction over the original source at the HLS level. The agent-driven phase delivered the remaining three orders of magnitude in less time. We do not claim this ratio generalises: the baseline phase included learning the tools, and the search phase benefited from the baseline's analysis of where the time went.

\textbf{What the agents did not do.} The rules stopped the agents from changing several things at once and from inflating evidence; they did not stop them from spending experiments on a direction that had stopped making sense, and the plan reviews are where that was corrected. All goal-level decisions were human: the choice of directions (Winograd and packing were explored because the course covered them), the precision freeze, the gate change, the scope decisions that closed open-ended workstreams (13 images are the test set; ROM parameter storage stays exploratory), and the choice of C22 over a rescue re-route of C36. The log also records clerical slips that the protocol caught after the fact rather than preventing: an experiment launched before its design existed (W26), a plan that specified three co-simulation calls where four were needed, and conditional implementation jobs that were never started. One of the costliest failures was a human instruction (running two implementations at once caused the W23 out-of-memory kill); the rule that replaced it came from the agent's own measurement of memory peaks.

\textbf{Threats to validity.} The quantitative findings are specific in three ways. The 2$\times$ HLS calibration is for one network, one device family and one tool version (Vitis HLS and Vivado 2026.1); it should be re-measured rather than assumed for a different design, although the direction of the error (pessimistic logic, optimistic clock, exact DSP count) is consistent with the design being dominated by wide adder trees that Vivado retimes. The bit-exact gate was checked on the 13 supplied images plus one Set5 image, and the 21-bit Conv2 accumulator relies on an observed range rather than a worst-case bound; a wider test set could in principle find an overflow that these images do not. And the experiment counts depend on the counting rule: 231 registered runs include retries, controls and verification runs of the same design at another evidence level; about 50 of the 71 design directories are distinct architectural points.

\textbf{Relation to other work.} Using LLMs to tune HLS designs is an active topic: reasoning-model agents that insert pragmas and restructure code with tool feedback \cite{ref21}, and, most directly comparable, an IBM study that drives Claude Code with Vitis HLS over benchmark kernels and reports mean speedups of about 7$\times$ \cite{ref22}. Those works optimise single kernels at the HLS-estimate level. This project differs in three ways: the object is a whole network carried through co-simulation, place-and-route and the board; the ranking deliberately refuses HLS estimates; and the contribution is not a search algorithm but a protocol that makes a general coding agent's search auditable by a reader who was not present. On the hardware side, Winograd on FPGAs \cite{ref18} and two-products-per-DSP packing are established techniques; what this report adds is the measured interaction between them on a small network, in particular that Winograd buys nothing until the fold is DSP-bound (Table~\ref{tab:22}, row 6 and W18) and that packing is only affordable where 8-bit weights are affordable (Table~\ref{tab:8}). We found the LLM-for-HLS literature after designing the method, and the comparison in this paragraph was not used to shape it.

\subsection{Limitations and future work}\label{sec:6.5}

Each limitation below has a concrete next step that is resource-neutral or nearly so, which matters because the device is 83 \% full on DSPs and 98 \% on CLBs.

\textbf{Memory access pattern.} The on-board cycle count is 1.8$\times$ the RTL count, and the AXI statistics attribute this to single-beat reads on the input and parameter paths, with the image channel waiting for data during 77.7 \% of the frame. The cause is in the HLS source: the loader reads, converts and writes to a stream in one pipelined loop with a division and a modulo in the address, so Vitis HLS cannot prove the accesses consecutive and generates one transaction per word \cite{ref12}. Section~\ref{sec:6.3} shows the fix at the HLS level: splitting the loader into a load-only row loop and a separate encode loop makes HLS infer a burst of length 65,025 on the image port with no change to the interface, the cycle estimate or the DSP count, at the cost of one BRAM18K. The remaining steps are to co-simulate that variant, package it, rebuild the system with the same HP0/HP1 topology and monitors, and compare 20 hardware-counted runs against the 489,584.5-cycle baseline, accepting it only if the output stays bit-exact and the measured average burst length exceeds 1. The parameter loaders need the same treatment, or, since the parameters are constant across frames, could stay on chip behind a \texttt{load\_params} switch as explored in Section~\ref{sec:3}. If the 216 k stall cycles were removed the board would approach the 1.638 ms the RTL predicts at the present clock. A more radical option is an AXI-Stream input fed by a DMA engine, which takes burst generation out of the accelerator entirely.

\textbf{Clock.} The IP was routed at 5 ns out of context, but the integrated system is implemented at the 5.333 ns Vivado derives from the PS preset, and on the board it runs at 166.665 MHz because the clock debugging stopped at a divider of 6 once the dependence of time on frequency had been established. The next divider, 5, gives 200 MHz, above the 187.5 MHz the system was implemented for; with 43 ps of setup slack at 5.333 ns it cannot be assumed to work. Reaching 200 MHz means re-implementing the full system at a 5 ns constraint, with the monitoring modules removed, and verifying timing and correctness again; the out-of-context result (+13 ps at 5 ns) suggests it is attainable but does not prove it. If it closes, the hardware time falls by a further 17 \% to about 2.45 ms with the present loader, or to the RTL's 1.365 ms once the memory path is fixed. Going below 5 ns is a different project: W36, C36 and C40 place the critical path inside the Winograd arithmetic with 8 to 11 logic levels, so a faster clock would need that engine re-pipelined.

\textbf{Resources.} With 82.7 \% of the DSPs and 97.7 \% of the CLBs in use there is no room for more buffers or parallel units; any architectural change beyond the two above would have to trade DSPs for LUTs or accept a different fold. The one place where a trade exists on paper is Conv1: F(4$\times$4, 3$\times$3) Winograd would need 324 DSPs instead of 576, but its transform constants are not powers of two, so the bit-exact-with-shifts property would be lost, and its window stream would exceed the 4,096-bit limit (Table~\ref{tab:7}). It was not pursued and is listed here as a known, unmeasured option.

\textbf{Software path and system throughput.} A full PYNQ invocation takes about 6 ms against 2.94 ms of hardware time (Table~\ref{tab:20}), so the system as delivered sustains about 167 frames per second while the accelerator alone would sustain 340. The difference is cache maintenance on four buffers, the NaN pre-fill, the output copy and the NaN/Inf check, all of which exist for verification. A production driver would keep the parameters resident and flush them once, drop the pre-fill and the check, and double-buffer input and output so that transfers overlap execution; the hardware time would then be the frame time. The 255$\times$255 frame size is fixed by the contract; a general accelerator would need variable sizes and a wider test set than the 14 images used here.

\textbf{Power.} Power was not measured. The only figure is Vivado's 5.8 W estimate at medium confidence (Section~\ref{sec:5}). The KV260 exposes its rails through PMBus \cite{ref11}, so a measurement under load, and the energy per frame that follows from it, would be a short addition and would allow a comparison with GPU and CPU implementations of the same network, which this report does not attempt.

\textbf{Verification scope of the method.} The search ranked designs on co-simulated cycles and routed clocks and never on the board, so the 2.15$\times$ gap of Section~\ref{sec:6.3} was invisible at every level the agents could reach. The most valuable methodological change is to add automated on-board measurement as the top evidence level of the protocol, so that memory-path effects feed back into the search instead of being discovered after it. The hardware for it already exists: the cycle monitor and the AXI statistics module of Section~\ref{sec:5} are exactly the instruments such a gate would read.
